Fig.~\ref{fig:schwarz_final_r_dd_2_x}, Fig.~\ref{fig:schwarz_final_r_dd_2_y} and Fig.~\ref{fig:schwarz_final_r_dd_2_z} are built exactly as their counterparts of the previous section, and the raw central panels are indeed identical: the uncoupled local solves do not depend on the transmission condition applied afterwards. Only the right panel changes, and it now reports the field obtained after the iteration \eqref{eq:local_update} has converged, rather than after a single application of \eqref{eq:local_solve}--\eqref{eq:truelocal}.

The iteration was run with the penalty $\rho = \sigma_{3D}/h_{\max} = 2.309\times10^{-3}$, stopping when the largest change between two consecutive iterates falls below $10^{-8}$. Convergence is monotone but slow in the tail: the reconstructed global error is $3.717\times10^{-1}$ after the first sweep, $1.268\times10^{-1}$ after five, $8.184\times10^{-2}$ after ten and $7.577\times10^{-2}$ after twenty, after which it is stationary to four digits. The first sweep carries no information between the boxes, since $u_{nbr}^{(0)} = 0$; the coupling only propagates from the second sweep onwards, and about twenty sweeps are needed for the interface traces to equilibrate across the eight subdomains.

At convergence the reconstructed global error is
\[
\frac{\lVert u^* - u^{\mathrm{schwarz}}\rVert_2}{\lVert u^*\rVert_2} = 7.553\times10^{-2},
\]
against the raw value $1.891\times10^{-1}$, i.e.\ a gain of a factor $2.5$, while the averaged local error falls from $3.415\times10^{-1}$ to $1.403\times10^{-1}$ (a factor $2.4$). As in the Robin case the two quantities improve by the same factor, confirming that the correction acts on the individual boxes and not through cancellation in the averaging. The improvement is again uniform: every subdomain gains between $2.1$ and $2.9$, and none deteriorates.

Measured on the cut plane, the relative $L^2$ error of the wall falls from $5.071\times10^{-1}$ to $2.066\times10^{-1}$ on $x=0$, from $3.760\times10^{-1}$ to $1.838\times10^{-1}$ on $y=0$, and from $2.684\times10^{-1}$ to $1.525\times10^{-1}$ on $z=0$, i.e.\ gains between $1.8$ and $2.5$.

Two differences with respect to the Robin-flux results deserve a comment. First, the gains are systematically smaller, by roughly a factor three on both the global and the wall errors. This is expected and is not a defect of the scheme: \eqref{eq:local_solve} evaluates the residual at $u^*$ and therefore consumes the exact solution, whereas \eqref{eq:local_update} only ever uses the neighbours' previous iterates. The Schwarz error is thus the one actually attainable when $u^*$ is unknown, and the Robin-flux figure of $2.625\times10^{-2}$ remains the lower bound that the surrogate is measured against.

Second, the error that survives the iteration is distributed differently. The median pointwise error on the wall improves by a factor between $3.6$ and $4.5$ (for instance from $1.893\times10^{-2}$ to $5.254\times10^{-3}$ on $x=0$), but the maximum barely moves, from $1.371\times10^{-1}$ to $1.234\times10^{-1}$ on $x=0$ and by a comparable $1.1$ on the other two planes. The residual error is therefore concentrated in a few isolated points, which the right panels locate on the boxes $(0,0,1)$ for the $x$ and $y$ cuts and $(0,1,1)$ for the $z$ cut, precisely where a vessel crosses the cut plane. This is the configuration the cross term of the previous section was introduced to handle: the Robin transmission condition acts on the $3$D trace alone and cannot represent a $1$D source whose averaging circle straddles the interface, and the Schwarz iteration, sharing the same $G_\gamma$, inherits the same blind spot.
